\documentclass[10pt,a4paper]{amsart}
\usepackage[margin=19mm]{geometry}
\usepackage{amsmath,amssymb,mathtools}
\usepackage[scr=rsfs,cal=euler]{mathalfa}
\usepackage{xfrac}
\usepackage[unicode,hidelinks,bookmarks=false]{hyperref}
\newcommand{\ie}{i.e.\ }
\newcommand{\cf}{cf.\ }
\newcommand{\resp}{respectively\ }
\newcommand{\Subsection}{\subsection}
\providecommand{\nobreakdash}{\nobreak\hbox{-}}
\setlength{\emergencystretch}{2em}
\multlinegap0pt
\usepackage{xcolor}
\usepackage{amssymb}
\newtheorem{theorem}{Theorem}
\title{Weaker quantization dimension results for self-similar measures}
\author{Saurabh Verma, Shivam Dubey}
\date{Source: arXiv:2601.20531v2 (CC BY 4.0)}
\begin{document}
\maketitle

\noindent\emph{Source and licence.} Weaker quantization dimension results for self-similar measures. Version arXiv:2601.20531v2 (CC BY 4.0), distributed by arXiv under the Creative Commons Attribution 4.0 International licence, http://creativecommons.org/licenses/by/4.0/, as recorded in the arXiv licence metadata for this exact version and shown on the abstract page https://arxiv.org/abs/2601.20531v2. Full licence notice: \texttt{licenses/arxiv-2601.20531-CC-BY-4.0.txt}.\par\smallskip

\noindent\textit{Editorial scope.} Counted unit: the article's theorem Lem24 (source0.tex lines 377--379). The article's complete original proof of it is delivered with it (source0.tex lines 380--398, one \texttt{proof} environment of 19 lines). No mathematics is written, repaired or re-derived: the statement and the proof are the article's own lines, in the article's own notation, and no retained line is substituted or re-flowed.  No further source-local context is needed: every cross-reference of every retained line resolves inside the two delivered spans.  Nothing was omitted from any retained span, so the package carries no omission record.  No retained line contains a citation, so the wrapper carries no bibliography and no bibliography entry.  No retained line contains a figure environment, an image-inclusion command or drawing code, so no image is delivered and the package declares no asset dependency.  Editorial formatting only, in addition: an AMS \texttt{amsart} preamble that restates the article's own declarations and macros used by the retained lines, namely the environments \texttt{theorem} on separate counters, together with the standard packages those lines need.  The retained environments print in this document as Theorem 1 in order of appearance, whereas the article prints the same items with the numbering of its own full document.  The complete native source of the article as distributed in the arXiv e-print for this exact version is bundled unchanged as \texttt{sources/193527/source0.tex}.
\begin{theorem}\label{Lem24}
    Let $\mu$ be a self-similar measure associated with a WIFS $(\{f_i\}_{i=1}^N; \\ \rho_1 ,\rho_2,\ldots,\rho_N)$ and $\mu_n$ be the self-similar measure associated with sub-WIFS $\mathcal{I}_n$ at $n$-th level. Then, $\mu_n$ is absolutely continuous with respect to $\mu$ i.e., $\mu_n << \mu.$
\end{theorem}

\begin{proof}
    Recall that, for a given probability vector $\{\varsigma_{\eta}= \frac{\rho_{ \eta \xi}}{ \sum_{\eta \in \Upsilon^n} \rho_{\eta \xi}}: \eta \in \Upsilon^n\},$ the map $\mathfrak{P}: \Omega(\mathbb{R}^m) \to \Omega(\mathbb{R}^m)$ defined by $\mathfrak{P}(\lambda)= \sum_{\eta \in \Upsilon^n} \varsigma_{\eta} \lambda \circ f_{\eta}^{-1}$ is contraction mapping with respect to the Monge-Kantorovich metric and by Banach contraction theorem it has a unique fixed point $\mu_n$. Also, independent of the choice of $\lambda \in \Omega(\mathbb{R}^m)$, the sequence $\mathfrak{P}^k(\lambda)$ converges to $\mu_n$ as $k$ approaches to infinity. Now, in order to prove that $\mu << \mu,$ first we show that if $\mu(A)=0$ for Borel set $A \subset \mathbb{R}^m,$ then $\mathfrak{P}^k(\mu)(A)= 0$ for all $k \in \mathbb{N}.$\\
    Consider, 
\begin{equation}\label{eq24}
\begin{aligned}
(\mathfrak{P} \circ \mathfrak{P})(\mu) = \mathfrak{P}\bigg(\sum_{\eta \in \Upsilon^n} \varsigma_{\eta} \mu \circ f_{\eta}^{-1} \bigg) = \sum_{\eta^1, \eta^2 \in \Upsilon^n}\varsigma_{\eta^1 \eta^2} \mu \circ f_{\eta^2 \eta^1}^{-1} 
\end{aligned}
\end{equation}
    Similarly, for $k \in \mathbb{N},$ we have 
    \[  \mathfrak{P}^k(\mu)= \sum_{\eta^1,\eta^2,\ldots,\eta^k \in \Upsilon^n} \varsigma_{\eta^1\eta^2\ldots\eta^k} \mu \circ f_{\eta^k \eta^{k-1}\ldots \eta^1}^{-1}.\]
    Since $\mu$ satisfies the self-similar equation $\mu= \sum_{i \in \Lambda} \rho_i \mu \circ f_{i}^{-1},$ we have 
    \[\mu= \sum_{i,j \in \Lambda} \rho_{ij} \mu \circ f_{ji}^{-1}.\]
    In $\zeta \in \Lambda^*$ is a word of length $l \in \mathbb{N}$ in the sub-IFS, then repeating the above process $l+n+k$ times we have 
    \begin{equation}\label{eq25}
    \begin{aligned}
        \mu(A) = \sum_{i^1,i^2,\ldots,i^{l+n+k} \in \Lambda} \rho_{i^1i^2\ldots i^{l+n+k}} \mu \circ f_{i^1i^2\ldots i^{l+n+k}}^{-1}(A)=0.
         \end{aligned}
    \end{equation}
    \end{proof}  

\end{document}
